\section{Introduction}
An LLM can produce plausible language while misclassifying social stereotypes,
changing its answer after a small input perturbation, or judging a morally relevant
action incorrectly. These failures matter when users depend on the model to interpret
people, documents, or requests. Part A evaluates one locally deployed open-weight
model using three subsets of TrustLLM: fairness, adversarial robustness, and machine
ethics. TrustLLM provides a common framework and task-specific datasets for studying
these dimensions~\cite{trustllm}.

The chosen tasks have reference labels, enabling automatic scoring without an external
LLM judge or hand-written keyword rules. Stereotype recognition tests the classification
of social associations; paired AdvGLUE examples test answer stability under adversarial
edits; ETHICS action-judgment examples test judgments about actions. Each is a limited
operational measure of a trust dimension. Recognizing a stereotype is different from
demonstrating equitable behavior in a real decision.

\section{Model and set-up}
\subsection{Local deployment and environment}
The evaluated endpoint was \texttt{google/gemma-4-e4b} in LM Studio 0.4.26+4 on an
Apple M4 computer with 16 GiB memory and macOS 27.0.1 (arm64). The server reported
architecture \texttt{gemma4}, MLX format, variant
\texttt{google/gemma-4-e4b@4bit}, four bits per weight, and a model file size of
6,861,915,950 bytes. The configured context window was 8,192 tokens. Google's E4B
model card describes 4.5 billion effective parameters and approximately 8 billion
including embeddings~\cite{gemma}; the local server's rounded \texttt{4B} metadata
should not be interpreted as its full parameter count. The exact upstream revision
and a hash of the local weights were not exposed by the recorded model metadata.

Inference used the local endpoint \texttt{http://127.0.0.1:1234/v1}, through the
OpenAI-compatible client protocol supported by LM Studio~\cite{lmstudio}. No remote
inference provider or additional moderation service was called. The model's own
instruction tuning remains part of the evaluated system. In Part A, each input was an independent
single-turn request, with temperature 0, maximum 192 output tokens,
\texttt{reasoning\_effort="none"}, and a 90-second timeout. Temperature 0 limits
sampling variation but does not establish bitwise reproducibility across runtimes.
No top-p or top-k override was supplied; server defaults were used.

Python 3.12.14 ran in a project-local \texttt{.venv} managed by uv~\cite{uv}.
The lockfile fixes the installed dependencies. Recorded packages were
\texttt{openai} 3.28.0, \texttt{huggingface-hub} 2.2.0, pandas 3.0.6,
scikit-learn 1.9.1, and Matplotlib 3.11.2~\cite{software,hf}.
The implementation separates download/sampling, prompts, inference persistence,
and offline scoring. The GitHub link and code availability details are in Appendix~\ref{app:code}.

\subsection{Data structure and sampling}
The three JSON files were downloaded through the authenticated Hugging Face Hub
client from \texttt{TrustLLM/TrustLLM-dataset}, pinned to revision
\nolinkurl{59fdb19b351609c93de79a9bf882af5801525d08}~\cite{hfdata}.
Each file has its own schema; none was forced into a shared dataset loader.
\begin{table}[H]
\centering\small
\begin{tabular}{p{2.4cm}p{4.0cm}p{4.5cm}r}
\toprule
Dimension & Subset & Reference label and metadata & Test inputs\\
\midrule
Fairness & \texttt{stereotype\_recognition} & Three classes; topic, sentence, context & 120\\
Robustness & \texttt{AdvGLUE} & Numeric task label; pair index, variant, attack method & 240\\
Machine ethics & \texttt{implicit\_ETHICS} & \texttt{wrong}/\texttt{not wrong}; input, topic & 120\\
\bottomrule
\end{tabular}
\caption{Selected TrustLLM subsets. Robustness comprises 120 complete pairs.}
\end{table}

The \texttt{prompt} field was sent verbatim as the user message. The \texttt{label}
field supplied the local reference answer and was never sent to the model.
Fairness \texttt{sentence}/\texttt{context} and ETHICS \texttt{input} retain the
underlying text for inspection; their \texttt{type}/\texttt{topic} fields describe
coverage. In AdvGLUE, \texttt{task} controls the label vocabulary;
\texttt{index} and \texttt{type} identify original/modified partners;
\texttt{method} and \texttt{data construction} describe perturbations. Full source
records are retained locally for auditing.

The downloaded pools contained 1,000 fairness rows, 912 AdvGLUE rows (456 complete
pairs), and 500 ETHICS rows (250 commonsense and 250 justice). The initial ETHICS
sample contained 59 commonsense and 61 justice examples. Python's
\texttt{random.Random(44)} shuffled each
classification label stratum or robustness task stratum. Before any smoke inference,
the procedure reserved a development pool of 30 fairness inputs, 20 ETHICS inputs,
and 12 robustness pairs; formal samples were drawn from the remaining rows in each
stratum. The smoke test used only 13 development inputs. Formal and development
sample IDs have zero overlap; robustness partners always stay in the same split.
No model training or prompt tuning on the formal set was performed. After the first
480-input run, a source audit identified incompatible task and label semantics in
the justice portion, as detailed in the ethics results. A second sample therefore
selected 120 commonsense-only ETHICS rows, balanced at 60 per label, with seed 44
and all original development IDs excluded. The selection did not depend on which
individual outputs were correct. Twenty-six rows overlap the first mixed test, but
all 120 were requested again in the documented refinement. Model settings, system
instructions, user prompts, and source labels were unchanged.

The primary reported sample contains 40 examples per fairness class, 60 per
commonsense ETHICS class,
and 30 AdvGLUE pairs for each of MNLI, QNLI, QQP, and SST-2. This yields 480 inputs
and 120 paired attack cases. The robustness set contains 211 distinct prompts,
satisfying at least 100 samples per dimension even after exact prompt deduplication.
The additional initial mixed-topic ETHICS run is retained as a diagnostic; the
Part A generated 600 responses in total, with 480 inputs contributing to the
primary three-dimension results. The ethics refinement was a disclosed post-run
measurement correction, rather than a preregistered blind test or a Part B mitigation.
Balancing keeps each class or task visible
within the available compute budget. These are stratified subset scores, rather
than estimates weighted to the full dataset's natural distribution. Fairness topics
were not stratified jointly: race 50, profession 46, gender 17, and religion 7.
Small topic strata limit group comparisons. Related classification contexts were
not explicitly grouped, so different row IDs do not guarantee semantic independence.
AdvGLUE reuses base examples for different attacks: the formal set has 91 distinct
original prompts (MNLI 23, QNLI 26, QQP 20, SST-2 22), and four exact prompts also
occur in the development pool. Thus the split separates attack-pair IDs, rather
than guaranteeing a holdout of all base texts. No prompts or settings were tuned
on these repeated texts. The pair-weighted robustness scores give more weight to
base texts represented by multiple attacks; they are not independent observations.
StereoSet underlies the stereotype task; AdvGLUE and ETHICS provide the other selected
material~\cite{stereoset,advglue,ethics}.

\subsection{Output protocol and metric definitions}
The shared system instruction requests English JSON with exactly two fields:
\texttt{label} and a short \texttt{explanation}. LM Studio receives a JSON schema
whose label enum contains the task's allowed classes plus \texttt{abstain}.
For fairness, the system maps the prompt's numbered options to stereotype,
anti-stereotype, and unrelated. For robustness the numeric references are decoded
as follows: SST-2 0/1 = negative/positive; QQP 0/1 = no/yes;
QNLI 0/1 = yes/no; MNLI 0/1/2 = yes/maybe/no. QNLI and QQP deliberately use different
positive-class encodings.

The original user prompt is unchanged, including its legacy output instruction,
but the higher-priority system instruction explicitly replaces that format with
JSON. The explanation is not used to infer the label. Parsing requires valid JSON,
the exact field names and types, an allowed label, and a completed
\texttt{finish\_reason="stop"}. Abstentions, invalid output, and truncation count
as incorrect in the full denominator. Technical failures block scoring. Only
connection errors, timeouts, HTTP 429, or server errors may be retried, at most twice;
semantic mistakes are never retried. This adapted protocol changes both format and
generation behavior, so it does not reproduce the official TrustLLM leaderboard.

For fairness, accuracy is the proportion matching the reference class, and macro-F1
is the unweighted mean of the three class F1 scores. For machine ethics, accuracy,
macro-F1, and balanced accuracy (the mean recall of \texttt{wrong} and
\texttt{not wrong}) are reported. Confusion matrices include abstention and invalid
prediction columns. For robustness, let $o_i$ and $m_i$ indicate correctness on
the original and modified inputs of pair $i$, with $N=120$:
\[
 A_{\mathrm{orig}}=\frac{\sum_i o_i}{N},\qquad
 A_{\mathrm{mod}}=\frac{\sum_i m_i}{N},\qquad
 \Delta=A_{\mathrm{orig}}-A_{\mathrm{mod}},\qquad
 \mathrm{ASR}=\frac{\sum_i o_i(1-m_i)}{\sum_i o_i}.
\]
Attack success rate (ASR) conditions on originally correct cases. If its denominator
is zero, it is undefined. Wrong-to-correct transitions are reported separately;
they can offset the aggregate accuracy drop. No combined trustworthiness score is
computed across these distinct tasks.
